% ============================================
% Assistant Professor Cover Letter – VCU Economics
% ============================================

\documentclass[11pt,letterpaper]{letter}

\usepackage{geometry}
\usepackage{microtype}
\usepackage[hidelinks]{hyperref}

\geometry{left=25mm,right=25mm,top=25mm,bottom=25mm}

\signature{Decory Edwards}
\address{
Department of Economics \\
Johns Hopkins University \\
Baltimore, MD 21218 \\
\texttt{dedwar65@jh.edu}
}

\begin{document}

\begin{letter}{Search Committee \\
Department of Economics \\
Virginia Commonwealth University \\
Richmond, VA \\ USA}

\opening{Dear Members of the Search Committee,}

I am writing to apply for the tenure-track Assistant Professor position in the Department of Economics at Virginia Commonwealth University. I am a Ph.D.\ candidate in Economics at Johns Hopkins University, advised by Professors Christopher D.~Carroll, Nicholas W.~Papageorge, and Stelios Fourakis, and I expect to receive my degree in May~2026. My research fields are macroeconomics, wealth and income dynamics, and computational economics.

My research agenda focuses on understanding the determinants of wealth inequality and the mechanisms through which heterogeneity in returns to wealth shapes aggregate outcomes. In my job-market paper, I estimate the degree of heterogeneity in individual portfolio returns using micro-data from the Survey of Consumer Finances and simulate a structural model in which households face idiosyncratic return risk. I show that allowing for persistent heterogeneity in returns substantially improves the model’s ability to match the empirical distribution of wealth, offering a tractable way to connect micro-level return dispersion to macroeconomic inequality.

In a second project, I use the Health and Retirement Study to examine the relationship between interpersonal trust and portfolio performance. Building on the literature that finds a hump-shaped relationship between trust and income, I show that a similar relationship holds for individual returns to wealth measured in the HRS data, highlighting a behavioral channel through which social attitudes shape saving and investment outcomes. This work also provides new estimates of the persistent component of returns to net worth, which motivate the calibration of heterogeneity in my structural model.

Virginia Commonwealth University’s Economics Department is an appealing environment for my work given its emphasis on applied, policy-relevant research and its strong engagement with undergraduate and graduate education in a diverse, urban public university setting. I am particularly interested in contributing to a department that values empirical rigor, quantitative methods, and research that speaks to real economic problems faced by households and policymakers. Richmond’s role as a state capital and regional economic hub also offers opportunities to connect economic research and teaching to policy and institutional contexts.

\textbf{Teaching.} I have experience teaching and supporting courses in \textit{Principles of Macroeconomics}, \textit{Intermediate Macroeconomic Theory}, and upper-level electives related to macroeconomic policy. My teaching emphasizes economic intuition, careful use of data, and clear links between theory and observed outcomes. At VCU, I would be well prepared to teach principles and intermediate macroeconomics, as well as contribute to field or applied courses aligned with the department’s curriculum and student interests. I value inclusive, student-centered instruction and am comfortable teaching students with diverse academic backgrounds.

I would be enthusiastic about building a long-term academic career at Virginia Commonwealth University and contributing to the department’s missions in research, teaching, and service.

Thank you very much for your time and consideration. I would welcome the opportunity to discuss my application further.

\closing{Sincerely,}

\end{letter}
\end{document}
